% Generated draft: requires mathematical and source review.
\documentclass[11pt]{article}
\usepackage[margin=25mm]{geometry}
\usepackage{fontspec}
\setmainfont{Latin Modern Roman}
\usepackage{amsmath,amssymb,mathtools,unicode-math}
\setmathfont{Latin Modern Math}
\usepackage{enumitem,needspace,xurl,hyperref}
\setlength{\parindent}{0pt}
\setlength{\parskip}{5pt}
\setlength{\emergencystretch}{3em}
\newcommand{\N}{\mathbb{N}}
\newcommand{\Z}{\mathbb{Z}}
\newcommand{\Q}{\mathbb{Q}}
\newcommand{\R}{\mathbb{R}}
\newcommand{\C}{\mathbb{C}}
\newcommand{\F}{\mathbb{F}}
\newcommand{\A}{\mathbb{A}}
\newcommand{\PP}{\mathbb{P}}
\newcommand{\E}{\mathbb{E}}
\newcommand{\eps}{\varepsilon}
\newcommand{\dd}{\,\mathrm{d}}
\begin{document}
\section*{UnsolvedMath ID: 700229}
\textbf{Source problem code:} catalogue entry
\subsection{Definitions and mathematical statement}
A group is \emph{soluble} if its derived series reaches the trivial subgroup after finitely many steps. A group is \emph{locally soluble} if every finitely generated subgroup is soluble, and is \emph{locally finite} if every finitely generated subgroup is finite.

For a group \(G\), let
\[
\mathcal M(G)=\{U\leq G: U\text{ is locally soluble and there is no locally soluble }W\text{ with }U<W\leq G\}.
\]
Here maximality is with respect to inclusion among locally soluble subgroups; a member of \(\mathcal M(G)\) is not required to be proper in \(G\). For \(g\in G\), write \(U^g=g^{-1}Ug\).

The source gives two alternative conditions, which we distinguish as
\[
\begin{aligned}
\mathfrak C_{\mathrm{conj}}
&=\{G: (\forall U,V\in\mathcal M(G))(\exists g\in G)\ U^g=V\},\\
\mathfrak C_{\mathrm{iso}}
&=\{G: (\forall U,V\in\mathcal M(G))\ U\cong V\}.
\end{aligned}
\]
In the second condition, isomorphism means abstract group isomorphism.

\textbf{Problem (R. Baer).} Describe the locally finite groups belonging to \(\mathfrak C_{\mathrm{conj}}\), or, under the source's alternative formulation, describe the locally finite groups belonging to \(\mathfrak C_{\mathrm{iso}}\).

The source states as background that a finite group belongs to its class \(\mathfrak C\) if and only if it is soluble. This is a supplied observation, not a result established in this draft. The parenthetical wording ``or at least isomorphic'' leaves the intended choice between the two conditions unresolved. Moreover, the question ``What can be said'' requests a description or structural information without specifying a particular conjectured characterization.
\subsection{Short English statement}
R. Baer's problem asks for a description of locally finite groups whose maximal locally soluble subgroups are all conjugate, or, in a weaker alternative, all abstractly isomorphic. The source does not select a single alternative or propose a specific characterization. Its finite-group observation and its open status are supplied source information.
\subsection{Sources}
\begin{itemize}
\item Local UnsolvedMath registry, permanent ID \(700229\), exact source problem code \texttt{catalogue entry}, title \emph{Kourovka Notebook Problem 4.19}. The supplied record attributes the problem to R. Baer, proposed in 1973, and labels it open.
\item \emph{The Kourovka Notebook (Unsolved Problems in Group Theory)}, original problem 4.19, 4th issue (1973); the supplied background identifies the 21st issue edition (2026), \texttt{arXiv:1401.0300v48}. This bibliographic attribution is supplied by the registry and has not been independently checked.
\item The supplied collection description names E. I. Khukhro and V. D. Mazurov as editors of the Kourovka Notebook collection covering Issues 1--21 (1965--2026).
\end{itemize}
\end{document}
